\chap{Chapter 17. IEC 62443 and Secure Architecture}
IEC 62443 is the industrial security standard everyone cites and few have read. This chapter gives its structure, what a security level really is, how plants are divided into zones and conduits, and the architecture patterns that follow, including what replaces the Purdue diagram.

\sect{17.1 Who each part is for}
\begin{tbl}{P{0.12\textwidth}P{0.3\textwidth}L}
Part & Written for & Subject\\\midrule
2-1 & Asset owner & The owner's security programme\\
2-4 & Service provider, integrator & Requirements on those who build and maintain systems\\
3-2 & Owner and integrator & Risk assessment; zones and conduits\\
3-3 & System integrator & System-level security requirements\\
4-1 & Product supplier & The secure development lifecycle\\
4-2 & Product supplier & Technical requirements for components\\
\end{tbl}
\prop{17.1}{Shared responsibility} Supplier, integrator and owner are in series. A certified component does not make a secure system, and a secure design does not survive a careless integrator, or a technician who leaves a default password.

\sect{17.2 Zones and conduits}
\defn{17.1}{Zone} A group of assets with the same security requirements.

\defn{17.2}{Conduit} The controlled communication path between two zones.

Zones are drawn from a risk assessment, not copied from the VLAN list: a VLAN is a network convenience, a zone is a statement about consequence. Safety systems normally get their own zone; so does wireless, and anything a vendor can reach remotely.

\law{17.1}{Possible conduits} Among z zones there are at most z(z − 1) / 2 pairs. Eight zones could be joined by 28 conduits; a good design uses only the few the process needs, each documented, filtered and monitored.

\sect{17.3 Security levels}
\begin{tbl}{P{0.12\textwidth}L}
Level & Protection against\\\midrule
SL 1 & Casual or accidental violation\\
SL 2 & Intentional attack with simple means and generic skills\\
SL 3 & Sophisticated means, moderate resources, industrial-control skills\\
SL 4 & Sophisticated means with extended resources: a state-level actor\\
\end{tbl}
\defn{17.3}{SL-T, SL-C and SL-A} The \emph{target} level a zone needs, from the risk assessment; the \emph{capability} a product can deliver when properly configured; and the level \emph{achieved} once installed. A product with SL-C 3 installed with default passwords has an SL-A nowhere near 3.

There are seven foundational requirements (FRs): FR1 identification and authentication, FR2 use control, FR3 system integrity, FR4 data confidentiality, FR5 restricted data flow, FR6 timely response to events, FR7 resource availability.

\law{17.2}{Security level is a vector} A security level is one number per FR. For each, the gap is g\_i = max(0, SL-T\_i − SL-A\_i), and a zone meets its target only when every gap is zero.

\wex{17.1} A treatment-plant control zone:
\begin{tbl}{P{0.18\textwidth}*{7}{>{\centering\arraybackslash}X}}
 & FR1 & FR2 & FR3 & FR4 & FR5 & FR6 & FR7\\\midrule
SL-T & 2 & 2 & 3 & 1 & 3 & 2 & 3\\
SL-A & 2 & 1 & 3 & 1 & 2 & 2 & 3\\
Gap & 0 & 1 & 0 & 0 & 1 & 0 & 0\\
\end{tbl}
Two gaps: operators share one login (FR2), and the conduit to the business network allows more than the historian traffic it exists for (FR5). Confidentiality, already at target, gets no money: availability and integrity normally outrank secrecy in a plant.

\defn{17.4}{Compensating countermeasure} When a legacy device cannot meet a requirement, the zone meets it instead: a firewall permitting only named hosts and functions, monitoring, physical access control, or a jump host with individual logins. This is how brownfield plants comply without replacing every controller.

\sect{17.4 Architecture after Purdue}
The Purdue levels (Section 1.4) began as a functional model and became a security model by accident. Cloud historians, cellular RTUs and sensors that talk straight to a vendor's platform have broken the neat diagram. A defensible position: Purdue as a rigid network diagram is dying; Purdue as a way of reasoning about trust boundaries is fine; zones and conduits are what you actually build.

\textbf{Zero trust, adapted for OT.} You cannot put an agent on a 15-year-old PLC, so trust is enforced \emph{around} devices: micro-segmentation at the cell level; remote access brokered through the DMZ with multi-factor authentication and session recording; protocol-aware firewalls that permit Modbus reads but block writes.

\law{17.3}{A dependency is a series element} If a local action needs a remote service, A\_action = A\_local × A\_remote. A local HMI at 0.9999 whose logins must reach a central identity server at 0.995 is usable only 0.9949 of the time: about 45 hours a year locked out instead of under one. Never let safety or basic operation depend on a remote authentication server; provide cached credentials or a break-glass account.

\defn{17.5}{DMZ rule} No traffic passes straight through the industrial DMZ. Sessions from either side terminate there, with different credentials on each side.

\defn{17.6}{Data diode} Hardware that physically permits traffic one way only, used to get data out of high-consequence zones. It cannot be misconfigured into letting traffic back in. Because no acknowledgement can return, diodes repeat each message: with loss p per copy and r copies, P\_lost = pʳ.

\textbf{The SANS Five ICS Critical Controls:} an ICS-specific incident response plan; a defensible architecture; ICS network visibility and monitoring; secure remote access; and risk-based vulnerability management.

\textbf{Cyber-Informed Engineering.} Assume the digital layer will be compromised and engineer out the worst outcome physically. A relief valve, a float switch hardwired to a pump stop, or a hardwired overspeed trip cannot be hacked.

\sect{17.5 Remember and ask}
\say{Security level is a vector, not a number; and SL-C is not SL-A.}
\say{What is the worst thing this process can physically do, and is there a non-digital barrier in front of it?}
\ask{What is the SL-T for this zone, and who signed off the risk assessment behind it?}
\ask{How does the remote access system behave when its cloud service is unreachable?}

\sect{Exercises}
\begin{enumerate}
\item Explain SL-T, SL-C and SL-A in one sentence each. (Definition 17.3)
\item Six zones. Find the maximum number of conduits. (15)
\item Target (3, 3, 3, 1, 2, 3, 3); achieved (2, 3, 3, 1, 2, 1, 3). Find the gap vector. ((1, 0, 0, 0, 0, 2, 0))
\item Name the five SANS ICS critical controls. (Section 17.4)
\item What does a data diode guarantee that a firewall does not? (Physically one-way traffic, whatever the configuration)
\end{enumerate}
